% AA-60
\chapter[Reprezentacje sl2 i układy korzeni]{Reprezentacje \texorpdfstring{$\mathfrak{sl}_2$}{sl2} i struktura układów korzeni}
\label{aa:ch:sl2-root-systems}

\providecommand{\atlasref}[2]{#2}


W rozdziale~\ref{aa:ch:cartan-root-decomposition} skonstruowaliśmy podalgebrę izomorficzną z $\mathfrak{sl}_2$ dla każdego korzenia. Teraz zbadamy jej reprezentacje i zastosujemy wyniki do reprezentacji dołączonej. W całym rozdziale $K$ jest algebraicznie domkniętym ciałem charakterystyki zero. Liczby całkowite utożsamiamy z ich obrazami w $K$.

\section{Reprezentacje nieprzywiedlne}
Niech $h,e,f$ będzie standardową bazą $\mathfrak{sl}_2(K)$, z relacjami $[h,e]=2e$, $[h,f]=-2f$, $[e,f]=h$. Operatory w reprezentacji oznaczamy przez $H,E,F$.

\begin{lemat}
\label{aa:lem:sl2-commutation}
Niech $K$ będzie algebraicznie domkniętym ciałem charakterystyki zero, a $\rho:\mathfrak{sl}_2(K)\to\mathfrak{gl}(V)$ reprezentacją. Oznaczmy $H=\rho(h)$, $E=\rho(e)$, $F=\rho(f)$, gdzie $h,e,f$ są standardowymi generatorami o relacjach podanych wyżej. Dla każdego całkowitego $j\ge1$ zachodzą
\[
 [H,F^j]=-2jF^j,\qquad [E,F^j]=jF^{j-1}(H-(j-1)I).
\]
\end{lemat}
\begin{proof}
Dla $j=1$ są to relacje definiujące. Z tożsamości $[A,BC]=[A,B]C+B[A,C]$ pierwszy wzór przechodzi z $j$ na $j+1$ przez dodanie $F^j[H,F]$. Dla drugiego wzoru otrzymujemy
\[
\begin{aligned}
 [E,F^{j+1}]&=jF^{j-1}(H-(j-1)I)F+F^jH\\
 &=jF^j(H-(j+1)I)+F^jH\\
 &=(j+1)F^j(H-jI),
\end{aligned}
\]
gdzie użyto $HF=FH-2F$.
\end{proof}

\begin{twierdzenie}
\label{aa:tw:sl2-simple-modules}
Niech $K$ będzie algebraicznie domkniętym ciałem charakterystyki zero. Każda niezerowa skończenie wymiarowa nieprzywiedlna reprezentacja $\mathfrak{sl}_2(K)$ jest izomorficzna z dokładnie jednym modułem $L(m)$, gdzie $m$ jest nieujemną liczbą całkowitą, o bazie $v_0,\ldots,v_m$ i działaniach
\[
 Hv_j=(m-2j)v_j,\quad Fv_j=v_{j+1},\quad
 Ev_j=j(m-j+1)v_{j-1},
\]
gdzie $v_{-1}=v_{m+1}=0$. Dla każdego $m\ge0$ wzory te definiują reprezentację nieprzywiedlną.
\end{twierdzenie}
\begin{proof}
Niech $V$ będzie taką reprezentacją. Skończone niepuste widmo $H$ zawiera $\lambda$ takie, że $\lambda+2$ nie należy do widma: w przeciwnym razie iterowane dodawanie $2$ dawałoby nieskończenie wiele wartości. Wybieramy $0\ne v$ z $Hv=\lambda v$. Relacja $HE=EH+2E$ daje $H(Ev)=(\lambda+2)Ev$, stąd $Ev=0$. Lemat~\ref{aa:lem:sl2-commutation} daje
\[
 HF^jv=(\lambda-2j)F^jv,\qquad
 EF^jv=j(\lambda-j+1)F^{j-1}v.
\]
Dla $j=0$ mamy $Hv=\lambda v$ i $Ev=0$; drugi wzór stosujemy dla $j\ge1$. Niezerowe $F^jv$ są liniowo niezależne jako wektory własne dla różnych wartości. Istnieje więc pierwszy indeks $m+1$ z $F^{m+1}v=0$. Stosując $E$, otrzymujemy $(m+1)(\lambda-m)F^mv=0$, czyli $\lambda=m$. Rozpiętość $v,Fv,\ldots,F^mv$ jest niezmiennicza względem wszystkich trzech operatorów, a więc równa $V$. Jednoznaczność $m$ wynika z $\dim V=m+1$.

Odwrotnie, w podanej bazie $[H,E]v_j=2Ev_j$, $[H,F]v_j=-2Fv_j$ oraz
$[E,F]v_j=((j+1)(m-j)-j(m-j+1))v_j=(m-2j)v_j$; wzory obejmują końce dzięki zerowym współczynnikom. Niezerowa podprzestrzeń niezmiennicza zawiera któryś $v_j$, ponieważ wielomianowe projekcje na przestrzenie własne $H$ ją zachowują. Kolejne zastosowania $E$ prowadzą do niezerowej wielokrotności $v_0$, gdyż współczynniki dla $1\le j\le m$ są niezerowe. Zastosowania $F$ dają całą bazę.
\end{proof}

\begin{przyklad}[Wielomiany jednorodne]
Na wielomianach jednorodnych stopnia $m$ w zmiennych $x,y$ kładziemy $H=x\partial_x-y\partial_y$, $E=x\partial_y$, $F=y\partial_x$. Dla $v_j=\frac{m!}{(m-j)!}x^{m-j}y^j$ otrzymujemy $Hv_j=(m-2j)v_j$, $Fv_j=v_{j+1}$ dla $j<m$, a $Fv_m=0$. Różniczkowanie po $y$ daje dla $j\ge1$ współczynnik $j(m-j+1)$ przy $v_{j-1}$, a $Ev_0=0$. Jest to model $L(m)$.
\end{przyklad}

\section{Pełna redukowalność}
\begin{lemat}[Operator Casimira]
\label{aa:lem:sl2-casimir}
Niech $K$ będzie algebraicznie domkniętym ciałem charakterystyki zero, a $H,E,F$ operatorami dowolnej reprezentacji $\mathfrak{sl}_2(K)$ odpowiadającymi $h,e,f$. Operator $C=H^2+2H+4FE$ komutuje z $H,E,F$. Na $L(m)$ działa przez $m(m+2)I$.
\end{lemat}
\begin{proof}
Mamy $[H,FE]=-2FE+2FE=0$ oraz
\[
\begin{aligned}
 [E,C]&=-2EH-2HE-4E+4HE=2[H,E]-4E=0,\\
 [F,C]&=2FH+2HF+4F-4FH=2[H,F]+4F=0.
\end{aligned}
\]
Na $v_0$ operator $C$ działa przez $m^2+2m$, bo $Ev_0=0$. Komutowanie z $F$ daje tę samą wartość na każdym $v_j=F^jv_0$.
\end{proof}

\begin{twierdzenie}
\label{aa:tw:sl2-complete-reducibility}
Niech $K$ będzie algebraicznie domkniętym ciałem charakterystyki zero. Każda skończenie wymiarowa reprezentacja $\mathfrak{sl}_2(K)$ w przestrzeni $V$ jest sumą prostą modułów $L(m)$, gdzie $m$ przebiega nieujemne liczby całkowite. Ich krotności są jednoznaczne.
\end{twierdzenie}
\begin{proof}
Indukujemy po wymiarze $V$, zaczynając od pustej sumy dla $V=0$. Wybieramy właściwy podmoduł $U$ największego wymiaru. Iloraz $V/U$ jest nieprzywiedlny, ponieważ przeciwobraz właściwego niezerowego podmodułu przeczyłby maksymalności. Zatem $V/U\simeq L(m)$, a z indukcji $U$ jest sumą prostą modułów $L(k)$.

Niech $q:V\to V/U$ będzie rzutem, $c=m(m+2)$, a $W$ uogólnioną przestrzenią własną operatora $C=H^2+2H+4FE$ dla $c$ (przyjmujemy przestrzeń zerową dla wartości spoza widma). Wszystkie uogólnione przestrzenie własne $C$ są podmodułami dzięki komutowaniu z $H,E,F$. Jeśli $(C-dI)^rx=0$ dla $d\ne c$, to $(c-d)^rq(x)=0$, więc $q(x)=0$. Z rozkładu na uogólnione przestrzenie własne wynika $q(W)=V/U$. Ponieważ $k(k+2)=m(m+2)$ dla nieujemnych liczb całkowitych zachodzi tylko dla $k=m$, podmoduł $U\cap W$ jest sumą kopii $L(m)$.

Widmo $H|_W$ zawiera się w $\{m,m-2,\ldots,-m\}$: w bazie dostosowanej do $U\cap W$ macierz jest blokowo górnotrójkątna, z blokami odpowiadającymi podmodułowi i ilorazowi. Wybieramy przeciwobraz $v$ wektora najwyższej wagi $v_0$ ilorazu w uogólnionej przestrzeni własnej $H|_W$ dla $m$. Taki przeciwobraz istnieje, ponieważ rozkład dowolnego przeciwobrazu na uogólnione przestrzenie własne ma obrazy zerowe poza składową dla $m$.

Relacje $(H-(\lambda+2)I)E=E(H-\lambda I)$ oraz $(H-(\lambda-2)I)F=F(H-\lambda I)$ pokazują po iteracji, że $E,F$ przesuwają uogólnione wagi odpowiednio o $2,-2$. Zatem $Ev=0$ oraz $F^{m+1}v=0$, bo wagi $m+2,-m-2$ nie występują. Wektor $u=(H-mI)v$ należy do $U\cap W$ i do uogólnionej przestrzeni własnej dla $m$. Na $U\cap W$ operator $H$ jest diagonalizowalny, więc $u$ jest zwykłym wektorem własnym lub zerem. Z Lematu~\ref{aa:lem:sl2-commutation} wynika
\[
 0=EF^{m+1}v=(m+1)F^mu.
\]
Na przestrzeni wagi $m$ w sumie kopii $L(m)$ operator $F^m$ jest iniektywny, gdyż na każdym składniku przekształca $v_0$ w $v_m$. Stąd $u=0$, czyli $Hv=mv$.

Wektory $v,Fv,\ldots,F^mv$ mają obrazy tworzące bazę ilorazu, więc są liniowo niezależne. Ich rozpiętość $L$ jest podmodułem dzięki wzorom komutacyjnym i $F^{m+1}v=0$. Rzut $q|_L$ jest izomorfizmem, zatem $V=U\oplus L$, gdzie $L\simeq L(m)$. To zamyka indukcję.

Jeśli $a_j$ jest wymiarem przestrzeni własnej $H$ dla $j\ge0$ (zero dla nieobecnej wagi), to każdy $L(k)$ wnosi jeden do $a_j$ dokładnie dla $k\ge j$ i $k-j$ parzystego. Zatem $a_j-a_{j+2}$ jest krotnością $L(j)$, niezależną od rozkładu.
\end{proof}

\section{Jednowymiarowość przestrzeni korzeniowych}
Niech odtąd $\mathfrak g$ będzie skończenie wymiarową półprostą algebrą Liego nad $K$, $\mathfrak h$ jej maksymalną podalgebrą toralną, a $\Phi$ zbiorem korzeni. Stosujemy oznaczenia $t_\alpha,h_\alpha,e_\alpha,f_\alpha$ z Twierdzenia~\ref{aa:tw:root-sl2-triple} oraz $\mathfrak g_0=\mathfrak h$, $\mathfrak g_\gamma=0$ dla niezerowych funkcjonałów niebędących korzeniami.

\begin{twierdzenie}
\label{aa:tw:root-dim-reduced}
Niech $\mathfrak g$ będzie skończenie wymiarową półprostą algebrą Liego nad algebraicznie domkniętym ciałem $K$ charakterystyki zero, $\mathfrak h$ jej maksymalną podalgebrą toralną, a $\Phi$ zbiorem korzeni względem $\mathfrak h$. Dla każdego $\alpha\in\Phi$ mamy $\dim\mathfrak g_\alpha=1$ oraz $K\alpha\cap\Phi=\{\alpha,-\alpha\}$. Dla wszystkich $\alpha,\beta\in\Phi$ liczba $\beta(h_\alpha)$ jest całkowita.
\end{twierdzenie}
\begin{proof}
Ograniczamy reprezentację dołączoną $\mathfrak g$ do podalgebry $\mathfrak{sl}_2$ związanej z $\alpha$. Z pełnej redukowalności i klasyfikacji jej modułów wszystkie wartości własne $\operatorname{ad}_{h_\alpha}$ są całkowite. Na $\mathfrak g_\beta$ wartością własną jest $\beta(h_\alpha)$.

Podmoduł $M=\mathfrak h\oplus\bigoplus_{k\in\mathbb Z\setminus\{0\}}\mathfrak g_{k\alpha}$ jest niezmienniczy na mocy relacji wagowych z AA-59. Jego wagi względem $h_\alpha$ są parzyste, równe $2k$, a przestrzenią wagi zero jest $\mathfrak h$. Rozkładamy $M$ na sumę prostą modułów nieprzywiedlnych. Ich najwyższe wagi są parzyste, bo występują w widmie $H|_M$. Każdy nietrywialny składnik nieprzywiedlny $L(2r)$ ma jednowymiarową przestrzeń wagi zero, na której $E$ działa niezerowo ze współczynnikiem $r(r+1)$. Składniki $L(0)$ są zerowane. Zatem liczba nietrywialnych składników równa się randze $\operatorname{ad}_{e_\alpha}|_{\mathfrak h}$, która wynosi jeden, ponieważ $[e_\alpha,t]=-\alpha(t)e_\alpha$.

Podalgebra związana z trójką jest podmodułem izomorficznym z $L(2)$: jej reprezentacja dołączona jest nieprzywiedlna dzięki prostocie $\mathfrak{sl}_2$, wykazanej w Przykładzie~\ref{aa:prz:lie-radical-examples}, i ma wymiar trzy. Rzuty tego podmodułu na składniki $L(0)$ są zerowe (niezerowy homomorfizm między nieprzywiedlnymi modułami jest izomorfizmem, gdyż jego jądro i obraz są podmodułami). Zatem jego rzut na jedyny nietrywialny składnik jest niezerowy i daje izomorfizm. Składnik jest więc $L(2)$. Wynika stąd $\dim\mathfrak g_\alpha=1$ i brak korzeni $k\alpha$ dla całkowitych $|k|\ge2$.

Jeśli $\beta=c\alpha\in\Phi$ dla $c\in K^\times$, to $t_\beta=ct_\alpha$ z jednoznaczności reprezentanta względem formy Killinga, a $h_\beta=c^{-1}h_\alpha$. Liczby $2c=\beta(h_\alpha)$ i $2/c=\alpha(h_\beta)$ są całkowite, a ich iloczyn to $4$. Stąd $c\in\{\pm1,\pm2,\pm1/2\}$. Przypadki $\pm2$ wykluczono, a $\pm1/2$ wykluczamy przez zastosowanie tego samego wyniku do $\beta$, gdyż wtedy $\alpha=\pm2\beta$. Korzenie $\alpha,-\alpha$ rzeczywiście należą do $\Phi$ na mocy Twierdzenia~\ref{aa:tw:root-killing-pairing}.
\end{proof}

\section{Ciągi korzeniowe}
\begin{twierdzenie}
\label{aa:tw:lie-root-string}
Niech $\mathfrak g$ będzie skończenie wymiarową półprostą algebrą Liego nad algebraicznie domkniętym ciałem $K$ charakterystyki zero, $\mathfrak h$ jej maksymalną podalgebrą toralną, a $\Phi$ zbiorem korzeni względem $\mathfrak h$. Dla $\alpha,\beta\in\Phi$ z $\beta\notin K\alpha$ istnieją jednoznaczne całkowite $p,q\ge0$ takie, że
\[
 \{k\in\mathbb Z:\beta+k\alpha\in\Phi\}=\{-p,\ldots,q\},
 \qquad \beta(h_\alpha)=p-q.
\]
Dla każdego $k\in\mathbb Z$, jeśli $\beta+k\alpha$ i $\beta+(k+1)\alpha$ są korzeniami, to $[e_\alpha,\mathfrak g_{\beta+k\alpha}]=\mathfrak g_{\beta+(k+1)\alpha}$.
\end{twierdzenie}
\begin{proof}
Suma $W=\bigoplus_{k\in\mathbb Z}\mathfrak g_{\beta+k\alpha}$ jest niezerowym podmodułem względem trójki związanej z $\alpha$, na mocy Lematu~\ref{aa:lem:toral-weight}. Żaden funkcjonał $\beta+k\alpha$ nie jest zerowy. Wagi operatora $\operatorname{ad}_{h_\alpha}|_W$ mają postać $b+2k$, gdzie $b=\beta(h_\alpha)$ jest całkowite; są parami różne, a każda występująca przestrzeń wagowa jest jednowymiarowa. Rozkładamy $W$ na moduły $L(m)$. Wszystkie ich najwyższe wagi mają tę samą parzystość co $b$. Dwa składniki miałyby wspólną wagę $0$ przy parzystym $b$ lub $1$ przy nieparzystym $b$, co przeczy jednowymiarowości. Stąd $W$ jest pojedynczym $L(m)$. Wagi tworzą ciąg od $-m$ do $m$ co dwa. Ponieważ $k=0$ występuje, indeksy tworzą przedział całkowity od $-p$ do $q$. Równania $b-2p=-m$, $b+2q=m$ dają $b=p-q$. Wzór na działanie $E$ mówi, że między sąsiednimi wagami działa niezerowo, poza najwyższą wagą, zatem jego obraz jest całą docelową prostą.
\end{proof}

\begin{wniosek}
\label{aa:cor:lie-root-reflection}
Dla $\alpha\in\Phi$ odwzorowanie $s_\alpha(\beta)=\beta-\beta(h_\alpha)\alpha$ na $\mathfrak h^*$ jest inwolucją zachowującą $\Phi$. Jeśli $\alpha,\beta,\alpha+\beta\in\Phi$, to $[\mathfrak g_\alpha,\mathfrak g_\beta]=\mathfrak g_{\alpha+\beta}$.
\end{wniosek}
\begin{proof}
Ponieważ $\alpha(h_\alpha)=2$, mamy $s_\alpha(\beta)(h_\alpha)=-\beta(h_\alpha)$, a więc $s_\alpha^2=I$. Dla korzeni proporcjonalnych do $\alpha$ odbicie zamienia $\alpha$ i $-\alpha$. W pozostałych przypadkach $-\beta(h_\alpha)=q-p$ leży w przedziale $[-p,q]$, więc $s_\alpha(\beta)\in\Phi$. Inwolucja daje równość $s_\alpha(\Phi)=\Phi$. Jeśli $\alpha+\beta\in\Phi$, to $\beta$ nie jest proporcjonalna do $\alpha$, ponieważ suma z $\pm\alpha$ daje zero lub $2\alpha$. Twierdzenie o ciągu korzeniowym dla $k=0$ i $\mathfrak g_\alpha=Ke_\alpha$ dają tezę o nawiasie.
\end{proof}

\section{Realizacja euklidesowa}

\noindent\emph{Powiązania tematyczne w Bibliotece matematycznej:} \atlasref{LA-12::ch:lin12-euklidesowe-unitarne}{LA-12: przestrzenie euklidesowe}.

\begin{twierdzenie}
\label{aa:tw:lie-root-realization}
Niech $\mathfrak g$ będzie niezerową skończenie wymiarową półprostą algebrą Liego nad algebraicznie domkniętym ciałem $K$ charakterystyki zero, $\mathfrak h$ jej maksymalną podalgebrą toralną, a $\Phi$ zbiorem korzeni względem $\mathfrak h$. Istnieje różnowartościowe odwzorowanie $\alpha\mapsto\bar\alpha$ zbioru $\Phi$ do rzeczywistej przestrzeni euklidesowej, którego obraz jest zredukowanym krystalograficznym układem korzeni w sensie Definicji~\ref{aa:def:root-system}, z relacją
\[
 \frac{2(\bar\beta,\bar\alpha)}{(\bar\alpha,\bar\alpha)}
 =\beta(h_\alpha)\qquad(\alpha,\beta\in\Phi).
\]
\end{twierdzenie}
\begin{proof}
Korzenie rozpinają $\mathfrak h^*$, a $\mathfrak h=0$ implikowałoby $\mathfrak g=0$ przez rozkład korzeniowy. Wybieramy korzenie $\alpha_1,\ldots,\alpha_r$ tworzące bazę $\mathfrak h^*$. Z niezdegenerowania formy Killinga na $\mathfrak h$ elementy $t_{\alpha_i}$ tworzą bazę $\mathfrak h$; przez przeskalowanie otrzymujemy bazę $h_i=h_{\alpha_i}$.

Niech $\kappa$ będzie formą Killinga $\mathfrak g$, a $B=(\kappa(h_i,h_j))_{i,j=1}^r$. Z rozkładu korzeniowego i jednowymiarowości przestrzeni korzeniowych mamy
\[
 B_{ij}=\sum_{\gamma\in\Phi}\gamma(h_i)\gamma(h_j).
\]
Wpisy są całkowite. Traktujemy $B$ jako macierz rzeczywistą. Dla $u\in\mathbb R^r$ zachodzi $u^{\mathsf T}Bu=\sum_\gamma(\sum_i u_i\gamma(h_i))^2$. Macierz $(\alpha_j(h_i))_{j,i}$ jest odwracalna nad $K$, bo opisuje parowanie dwóch baz. Ma całkowite wpisy, więc jej wyznacznik jest niezerową liczbą całkowitą i jest odwracalna nad $\mathbb R$. Dla $u\ne0$ co najmniej jeden składnik sumy kwadratów jest dodatni. Zatem $B$ jest dodatnio określona.

Na $V=\mathbb R^r$ kładziemy $(u,v)=u^{\mathsf T}B^{-1}v$. Macierz $B^{-1}$ jest symetryczna; dla $u\ne0$ wektor $w=B^{-1}u$ jest niezerowy i $u^{\mathsf T}B^{-1}u=w^{\mathsf T}Bw>0$. Jest to więc iloczyn skalarny. Definiujemy $\bar\gamma=(\gamma(h_i))_{i=1}^r$ jako rzeczywisty wektor kolumnowy. Różnowartościowość i niezerowość wynikają z tego, że $h_i$ jest bazą; wektory $\bar\alpha_i$ rozpinają $V$ dzięki powyższemu wyznacznikowi. Współrzędne $t_\alpha$ w bazie $h_i$ to $B^{-1}\bar\alpha$, ponieważ $\kappa(t_\alpha,h_j)=\alpha(h_j)$. Są wymierne, a rachunek na nich w $K$ zgadza się z rachunkiem rzeczywistym. Stąd $\beta(t_\alpha)=\bar\beta^{\mathsf T}B^{-1}\bar\alpha$. Wstawiając $h_\alpha=2t_\alpha/\alpha(t_\alpha)$, otrzymujemy żądany iloraz.

Odbicie euklidesowe względem hiperpłaszczyzny prostopadłej do $\bar\alpha$ odpowiada zatem $s_\alpha$, więc zachowuje obraz $\Phi$. Liczby Cartana są całkowite. Jeśli dwa wektory korzeni są proporcjonalne nad $\mathbb R$, współczynnik proporcjonalności jest wymierny jako iloraz dwóch niezerowych współrzędnych całkowitych. Korzenie są więc proporcjonalne również nad $K$, a współczynnik musi być równy $\pm1$. Wykazano wszystkie aksjomaty.
\end{proof}

\begin{przyklad}[Ciąg korzeniowy w $\mathfrak{sl}_3$]
W oznaczeniach Przykładu~\ref{aa:prz:sl3-root-decomposition} bierzemy $\alpha=\alpha_{12}$, $\beta=\alpha_{23}$. Mamy $\beta(h_\alpha)=-1$, $\beta+\alpha=\alpha_{13}$, natomiast $\beta-\alpha$ i $\beta+2\alpha$ nie są korzeniami. Zatem $p=0,q=1$. Nawiasy $[E_{12},E_{23}]=E_{13}$ i $[E_{21},E_{13}]=E_{23}$ pokazują, że te dwie przestrzenie tworzą $L(1)$.
\begin{center}
\begin{tikzpicture}[every node/.style={font=\small}]
\node (a) at (-2,0) {$\mathfrak g_\beta$};
\node (b) at (2,0) {$\mathfrak g_{\beta+\alpha}$};
\draw[-{Stealth},bend left=20] (a) to node[above] {$\operatorname{ad}_{e_\alpha}$} (b);
\draw[-{Stealth},bend left=20] (b) to node[below] {$\operatorname{ad}_{f_\alpha}$} (a);
\node at (-2,-.9) {waga $-1$};
\node at (2,-.9) {waga $1$};
\end{tikzpicture}
\end{center}
Odbicie $s_\alpha$ zamienia $\beta$ z $\beta+\alpha$.
\end{przyklad}

Teoria z rozdziału~\ref{aa:ch:root-systems} stosuje się do otrzymanego układu korzeni. Klasyfikacja układów korzeni nie jest jeszcze klasyfikacją algebr Liego: potrzebne są dowody odtworzenia algebry z danych korzeniowych i istnienia odpowiednich algebr.
